\subsection{Summary of Findings}

We present an execution-aware pathfinding framework for stablecoin cross-exchange arbitrage using heuristic-guided A* search. Our approach bridges the gap between theoretical arbitrage detection and practical execution by incorporating real-world constraints directly into the search process.

\textbf{Key Contributions:}

\begin{enumerate}
    \item \textbf{Novel Algorithmic Framework}: We introduce A* search with domain-specific heuristics (H1--H4) that incorporate liquidity, slippage, risk, and parallel exploration, providing a practical alternative to theoretical detection algorithms.

    \item \textbf{Performance Improvements}: H4 achieves 10--20× faster execution (4.7s vs.\ 66--105s) while maintaining competitive profitability, demonstrating that risk-aware heuristics improve efficiency without sacrificing solution quality.

    \item \textbf{High Reliability}: H1 and H3 achieve 100\% success rates in Monte Carlo simulations, while H2 and H4 achieve 80\% with superior speed-efficiency trade-offs.

    \item \textbf{Scalable Profitability}: Profit percentages remain constant (0.42\%) across order sizes from \$1,000 to \$100,000, indicating accurate fee modeling and scalable capital deployment.

    \item \textbf{Optimal Path Structure}: Successful paths consistently require 3 hops, revealing optimal arbitrage structure and enabling efficient search space pruning.

    \item \textbf{Baseline Validation}: Bellman-Ford negative cycle detection fails in our execution-aware model, validating that theoretical algorithms cannot account for real-world constraints.
\end{enumerate}

Our work demonstrates that execution constraints can be effectively incorporated into heuristic functions, enabling real-time evaluation during pathfinding and bridging the gap between theoretical detection and practical execution systems.

\subsection{Practical Implications}

Our framework provides actionable guidance: H4 is recommended for production deployment requiring speed (\$38.45/sec efficiency), H1 for conservative strategies prioritizing reliability (100\% success rate), H2 for large capital deployments where slippage dominates, and H3 for opportunity discovery (highest average profit: \$203.49). Linear profit scaling suggests the system scales to larger capital deployments, with sub-second to minute-scale execution times sufficient for real-time deployment. Path structure analysis reveals that optimal arbitrage requires simple 3-hop paths, enabling efficient search space exploration.

\subsection{Future Directions}

Promising directions include hybrid heuristics with adaptive weighting, machine learning for dynamic parameter optimization, real-time execution integration, and multi-objective optimization. Large-scale Monte Carlo simulations and historical backtesting would provide additional validation. Expanding to additional exchanges, cross-asset arbitrage, and decentralized exchange integration would demonstrate broader applicability.

Our research establishes execution-aware pathfinding as a practical and effective approach to stablecoin arbitrage, with clear pathways for continued development and production deployment.



